% Chunk: \S5.9 Massless Particle Fields
% Source: Weinberg QFT Vol. I, physical pp. 269--278 (printed pp. 246--255); begins at the \S5.9 heading and ends immediately before the Problems heading.
% Equations: (5.9.1)--(5.9.41). Figures: none. Footnotes: 2.
% Uncertainties: none.
% Status: source-reviewed, compile-clean, and visually verified; frozen for primary-agent integration.

\subsection{Massless Particle Fields}

Up to this point we have dealt only with the fields of massive particles. For
some of these fields, such as the scalar and Dirac fields discussed in
Sections 5.2 and 5.5, there is no special problem in passing to the limit of
zero mass. On the other hand, we saw in Section 5.3 that there is a difficulty
in taking the zero-mass limit of the vector field for a particle of spin one:
at least one of the polarization vectors blows up in this limit. In fact, we
shall see in this section that the creation and annihilation operators for
physical massless particles of spin $j\geq1$ cannot be used to construct all of
the irreducible $(A,B)$ fields that can be constructed for finite mass. This
peculiar limitation on field types will lead us naturally to the introduction
of gauge invariance.

Just as we did for massive particles, let us attempt to construct a general
free field for a massless particle as a linear combination of the annihilation
operators $a_{\mathbf p,\sigma}$ for particles of momentum $\mathbf p$ and
helicity $\sigma$, and the corresponding creation operators
$a^{c\dagger}_{\mathbf p,\sigma}$ for the antiparticles:\footnote{We deal here
with only a single species of particle, and drop the species label $n$. Also,
$\kappa$ and $\lambda$ are constant coefficients to be determined by the
requirement of causality with some convenient choice of normalization of the
coefficient functions $u_\ell$ and $v_\ell$.}
\begin{align}
\psi_\ell(x)={}&(2\pi)^{-3/2}\int d^3p\sum_\sigma
\bigl[\kappa a_{\mathbf p,\sigma}u_\ell(\mathbf p,\sigma)e^{ip\cdot x}
+\lambda a^{c\dagger}_{\mathbf p,\sigma}
v_\ell(\mathbf p,\sigma)e^{-ip\cdot x}\bigr]
\tag{5.9.1}\label{eq:5.9.1}
\end{align}
where now $p^0\equiv|\mathbf p|$. The creation operators transform just like
the one-particle states in Eq.~\eqref{eq:2.5.42}
\begin{equation}
U(\Lambda)a^\dagger_{\mathbf p,\sigma}U^{-1}(\Lambda)=
\sqrt{\frac{(\Lambda p)^0}{p^0}}\exp(i\sigma\theta(p,\Lambda))
a^\dagger_{\mathbf p_\Lambda,\sigma},
\tag{5.9.2}\label{eq:5.9.2}
\end{equation}
\begin{equation}
U(\Lambda)a^{c\dagger}_{\mathbf p,\sigma}U^{-1}(\Lambda)=
\sqrt{\frac{(\Lambda p)^0}{p^0}}\exp(i\sigma\theta(p,\Lambda))
a^{c\dagger}_{\mathbf p_\Lambda,\sigma},
\tag{5.9.3}\label{eq:5.9.3}
\end{equation}
and hence also
\begin{equation}
U(\Lambda)a_{\mathbf p,\sigma}U^{-1}(\Lambda)=
\sqrt{\frac{(\Lambda p)^0}{p^0}}\exp(-i\sigma\theta(p,\Lambda))
a_{\mathbf p_\Lambda,\sigma},
\tag{5.9.4}\label{eq:5.9.4}
\end{equation}
where $p_\Lambda\equiv\Lambda p$, and $\theta$ is the angle defined by Eqs.
\eqref{eq:2.5.43}. Hence if we want the field to transform according to some
representation $D(\Lambda)$ of the homogeneous Lorentz group
\begin{equation}
U(\Lambda)\psi_\ell(x)U^{-1}(\Lambda)=
\sum_{\bar\ell}D_{\ell\bar\ell}(\Lambda^{-1})
\psi_{\bar\ell}(\Lambda x),
\tag{5.9.5}\label{eq:5.9.5}
\end{equation}
then we must take the coefficient functions $u$ and $v$ to satisfy the
relations
\begin{equation}
u_{\bar\ell}(\mathbf p_\Lambda,\sigma)
\exp(i\sigma\theta(p,\Lambda))=
\sqrt{\frac{p^0}{(\Lambda p)^0}}
\sum_\ell D_{\bar\ell\ell}(\Lambda)u_\ell(\mathbf p,\sigma),
\tag{5.9.6}\label{eq:5.9.6}
\end{equation}
\begin{equation}
v_{\bar\ell}(\mathbf p_\Lambda,\sigma)
\exp(-i\sigma\theta(p,\Lambda))=
\sqrt{\frac{p^0}{(\Lambda p)^0}}
\sum_\ell D_{\bar\ell\ell}(\Lambda)v_\ell(\mathbf p,\sigma)
\tag{5.9.7}\label{eq:5.9.7}
\end{equation}
in place of Eqs.~\eqref{eq:5.1.19} and \eqref{eq:5.1.20}. (Again,
$p_\Lambda\equiv\Lambda p$.) As in the massive particle case, we can satisfy
these requirements by setting (in place of Eqs.~\eqref{eq:5.1.21} and \eqref{eq:5.1.22})
\begin{align}
u_{\bar\ell}(\mathbf p,\sigma)&=
\sqrt{\frac{|\mathbf k|}{p^0}}\sum_\ell
D_{\bar\ell\ell}(\mathcal L(p))u_\ell(\mathbf k,\sigma),
\tag{5.9.8}\label{eq:5.9.8}\\
v_{\bar\ell}(\mathbf p,\sigma)&=
\sqrt{\frac{|\mathbf k|}{p^0}}\sum_\ell
D_{\bar\ell\ell}(\mathcal L(p))v_\ell(\mathbf k,\sigma),
\tag{5.9.9}\label{eq:5.9.9}
\end{align}
where $\mathbf k$ is a standard momentum, say $(0,0,k)$, and
$\mathcal L(p)$ is a standard Lorentz transformation that takes a massless
particle from momentum $\mathbf k$ to momentum $\mathbf p$. Also, in place of
Eqs.~\eqref{eq:5.1.23} and \eqref{eq:5.1.24}, the coefficient functions at the standard momentum
must satisfy
\begin{align}
u_{\bar\ell}(\mathbf k,\sigma)\exp(i\sigma\theta(k,W))
&=\sum_\ell D_{\bar\ell\ell}(W)u_\ell(\mathbf k,\sigma),
\tag{5.9.10}\label{eq:5.9.10}\\
v_{\bar\ell}(\mathbf k,\sigma)\exp(-i\sigma\theta(k,W))
&=\sum_\ell D_{\bar\ell\ell}(W)v_\ell(\mathbf k,\sigma)
\tag{5.9.11}\label{eq:5.9.11}
\end{align}
where $W^\mu{}_\nu$ is an arbitrary element of the `little group' for
four-momentum $k=(\mathbf k,|\mathbf k|)$, i.e., an arbitrary Lorentz
transformation that leaves this four-momentum invariant.

We can extract the content of Eqs.~\eqref{eq:5.9.10} and \eqref{eq:5.9.11} by considering
separately the two kinds of little-group elements in Eq.~\eqref{eq:2.5.28}. For a
rotation $R(\theta)$ by an angle $\theta$ around the $z$-axis, given by Eq.
\eqref{eq:2.5.27},
\[
R^\mu{}_\nu(\theta)=
\begin{bmatrix}
\cos\theta&-\sin\theta&0&0\\
\sin\theta&\cos\theta&0&0\\
0&0&1&0\\
0&0&0&1
\end{bmatrix},
\]
we find from Eqs.~\eqref{eq:5.9.10} and \eqref{eq:5.9.11}
\begin{align}
u_{\bar\ell}(\mathbf k,\sigma)e^{i\sigma\theta}
&=\sum_\ell D_{\bar\ell\ell}(R(\theta))u_\ell(\mathbf k,\sigma),
\tag{5.9.12}\label{eq:5.9.12}\\
v_{\bar\ell}(\mathbf k,\sigma)e^{-i\sigma\theta}
&=\sum_\ell D_{\bar\ell\ell}(R(\theta))v_\ell(\mathbf k,\sigma).
\tag{5.9.13}\label{eq:5.9.13}
\end{align}
For combined rotations and boosts $S(\alpha,\beta)$ in the $x-y$ plane, given
by \eqref{eq:2.5.26},
\[
S^\mu{}_\nu(\alpha,\beta)=
\begin{bmatrix}
1&0&-\alpha&\alpha\\
0&1&-\beta&\beta\\
\alpha&\beta&1-\gamma&\gamma\\
\alpha&\beta&-\gamma&1+\gamma
\end{bmatrix},
\qquad \gamma\equiv(\alpha^2+\beta^2)/2,
\]
Eqs.~\eqref{eq:5.9.10} and \eqref{eq:5.9.11} give
\begin{align}
u_{\bar\ell}(\mathbf k,\sigma)&=
\sum_\ell D_{\bar\ell\ell}(S(\alpha,\beta))u_\ell(\mathbf k,\sigma),
\tag{5.9.14}\label{eq:5.9.14}\\
v_{\bar\ell}(\mathbf k,\sigma)&=
\sum_\ell D_{\bar\ell\ell}(S(\alpha,\beta))v_\ell(\mathbf k,\sigma).
\tag{5.9.15}\label{eq:5.9.15}
\end{align}
Eqs.~\eqref{eq:5.9.12}--\eqref{eq:5.9.15} are the conditions that determine the coefficient
functions $u$ and $v$ at the standard momentum $\mathbf k$; Eqs.~\eqref{eq:5.9.8} and
\eqref{eq:5.9.9} then give them at arbitrary momenta. The equations for $v$ are just the
complex conjugates of the equations for $u$, so with a suitable adjustment of
the constants $\kappa$ and $\lambda$ we may normalize the coefficient
functions so that
\begin{equation}
v_\ell(\mathbf p,\sigma)=u_\ell(\mathbf p,\sigma)^*.
\tag{5.9.16}\label{eq:5.9.16}
\end{equation}

The problem is that we cannot find a $u_\ell$ that satisfies Eq.~\eqref{eq:5.9.14} for
general representations of the homogeneous Lorentz group, even for those
representations for which it is possible to construct fields for particles of
a given helicity in the case $m\neq0$.

To see what goes wrong here, let's try to construct the four-vector
$[(\tfrac12,\tfrac12)]$ field for a massless particle of helicity $\pm1$. In
the four-vector representation, we have simply
\[
D^\mu{}_\nu(\Lambda)=\Lambda^\mu{}_\nu.
\]
It is conventional to write the coefficient function $u_\mu$ here in terms of
a `polarization vector' $e_\mu$:
\begin{equation}
u_\mu(\mathbf p,\sigma)\equiv(2p^0)^{-1/2}e_\mu(\mathbf p,\sigma),
\tag{5.9.17}\label{eq:5.9.17}
\end{equation}
so that Eq.~\eqref{eq:5.9.8} gives
\begin{equation}
e^\mu(\mathbf p,\sigma)=\mathcal L(p)^\mu{}_\nu e^\nu(\mathbf k,\sigma).
\tag{5.9.18}\label{eq:5.9.18}
\end{equation}
Also, Eqs.~\eqref{eq:5.9.12} and \eqref{eq:5.9.14} read here
\begin{align}
e^\mu(\mathbf k,\sigma)e^{i\sigma\theta}
&=R(\theta)^\mu{}_\nu e^\nu(\mathbf k,\sigma),
\tag{5.9.19}\label{eq:5.9.19}\\
e^\mu(\mathbf k,\sigma)&=S(\alpha,\beta)^\mu{}_\nu
e^\nu(\mathbf k,\sigma).
\tag{5.9.20}\label{eq:5.9.20}
\end{align}
Eq.~\eqref{eq:5.9.19} requires that (up to a constant which can be absorbed into the
coefficients $\kappa$ and $\lambda$),
\begin{equation}
e^\mu(\mathbf k,\pm1)=(1,\pm i,0,0)/\sqrt2.
\tag{5.9.21}\label{eq:5.9.21}
\end{equation}
But then Eq.~\eqref{eq:5.9.20} would require also that $\alpha+i\beta=0$, which is
impossible for general real $\alpha,\beta$. We therefore cannot satisfy the
fundamental requirement \eqref{eq:5.9.14} or \eqref{eq:5.9.10}; instead, we have here
\begin{align}
D^\mu{}_\nu(W(\theta,\alpha,\beta))e^\nu(\mathbf k,\pm1)
&=S^\mu{}_\lambda(\alpha,\beta)R^\lambda{}_\nu(\theta)
e^\nu(\mathbf k,\pm1)
\notag\\
&=\exp(\pm i\theta)\left\{e^\mu(\mathbf k,\pm1)
+\frac{\alpha\pm i\beta}{\sqrt{2}|\mathbf k|}k^\mu\right\}.
\tag{5.9.22}\label{eq:5.9.22}
\end{align}
We have thus come to the conclusion that no four-vector field can be
constructed from the annihilation and creation operators for a particle of
mass zero and helicity $\pm1$.

Let's temporarily close our eyes to this difficulty, and go ahead anyway,
using Eqs.~\eqref{eq:5.9.18} and \eqref{eq:5.9.21} to define a polarization vector for arbitrary
momentum, and take the field as
\begin{align}
a_\mu(x)={}&\int \frac{d^3 p}{\sqrt{(2\pi)^3 2p^0}} \,
\cdot\sum_{\sigma=\pm1}\bigl[e_\mu(\mathbf p,\sigma)e^{ip\cdot x}
a_{\mathbf p,\sigma}+e_\mu(\mathbf p,\sigma)^*e^{-ip\cdot x}
a^{c\dagger}_{\mathbf p,-\sigma}\bigr].
\tag{5.9.23}\label{eq:5.9.23}
\end{align}
We will come back later to consider how such a field can be used as an
ingredient in a physical theory.

The field \eqref{eq:5.9.23} of course satisfies
\begin{equation}
\Box a^\mu(x)=0.
\tag{5.9.24}\label{eq:5.9.24}
\end{equation}
Other properties of the field follow from those of the polarization vector.
(We shall need these properties of the polarization vector later when we come
to quantum electrodynamics.) Note that the Lorentz transformation $\mathcal
L(p)$ that takes a massless particle momentum from $\mathbf k$ to $\mathbf p$
may be written as a `boost' $\mathcal B(|\mathbf p|)$ along the $z$-axis which
takes the particle from energy $|\mathbf k|$ to energy $|\mathbf p|$, followed
by a standardized rotation $R(\widehat{\mathbf p})$ that takes the $z$-axis
into the direction of $\mathbf p$. Since $e^\nu(\mathbf k,\pm1)$ is a purely
spatial vector with only $x$ and $y$ components, it is unaffected by the boost
along the $z$-axis, and so
\begin{equation}
e^\mu(\mathbf p,\pm1)=R(\widehat{\mathbf p})^\mu{}_\nu
e^\nu(\mathbf k,\pm1).
\tag{5.9.25}\label{eq:5.9.25}
\end{equation}
In particular, $e^0(\mathbf k,\pm1)=0$ and
$\mathbf k\cdot\mathbf e(\mathbf k,\pm1)=0$ so
\begin{align}
e^0(\mathbf p,\pm1)&=0
\tag{5.9.26}\label{eq:5.9.26}\\
\intertext{and}
\mathbf p\cdot\mathbf e(\mathbf p,\pm1)&=0.
\tag{5.9.27}\label{eq:5.9.27}
\end{align}
It follows that
\begin{align}
a^0(x)&=0
\tag{5.9.28}\label{eq:5.9.28}\\
\intertext{and}
\boldsymbol\nabla\cdot\mathbf a(x)&=0.
\tag{5.9.29}\label{eq:5.9.29}
\end{align}

As we shall see in Chapter 9, these are the conditions satisfied by the vacuum
vector potential of electrodynamics in what is called Coulomb or radiation
gauge.

The fact that $a^0$ vanishes in all Lorentz frames shows vividly that $a^\mu$
cannot be a four-vector. Instead, Eq.~\eqref{eq:5.9.22} shows that for a general momentum
$\mathbf p$ and a general Lorentz transformation $\Lambda$, in place of Eq.
\eqref{eq:5.9.6} we have
\begin{equation}
e^\mu(\mathbf p_\Lambda,\pm1)\exp(\pm i\theta(p,\Lambda))
=D^\mu{}_\nu(\Lambda)e^\nu(\mathbf p,\pm1)
+p^\mu\Omega_\pm(p,\Lambda),
\tag{5.9.30}\label{eq:5.9.30}
\end{equation}
so that under a general Lorentz transformation
\begin{equation}
U(\Lambda)a_\mu(x)U^{-1}(\Lambda)=
\Lambda^\nu{}_\mu a_\nu(\Lambda x)+\partial_\mu\Omega(x,\Lambda),
\tag{5.9.31}\label{eq:5.9.31}
\end{equation}
where $\Omega(x,\Lambda)$ is a linear combination of annihilation and creation
operators, whose precise form will not concern us here. As we will see in more
detail in Chapter 8, we will be able to use a field like $a^\mu(x)$ as an
ingredient in Lorentz-invariant physical theories if the couplings of
$a^\mu(x)$ are not only formally Lorentz-invariant (that is, invariant under
formal Lorentz transformations under which $a^\mu\to\Lambda^\mu{}_\nu
a^\nu$), but are also invariant under the `gauge' transformations
$a_\mu\to a_\mu+\partial_\mu\Omega$. This is accomplished by taking the
couplings of $a_\mu$ to be of the form $a_\mu j^\mu$, where $j^\mu$ is a
four-vector current with $\partial_\mu j^\mu=0$.

Although there is no ordinary four-vector field for massless particles of
helicity $\pm1$, there is no problem in constructing an antisymmetric tensor
field for such particles. From Eq.~\eqref{eq:5.9.22} and the invariance of $k^\mu$ under
the little group we see immediately that
\begin{align}
&D^\mu{}_\rho(W(\theta,\alpha,\beta))
D^\nu{}_\sigma(W(\theta,\alpha,\beta))
\bigl(k^\rho e^\sigma(\mathbf k,\pm1)-k^\sigma e^\rho(\mathbf k,\pm1)\bigr)
\notag\\
&\qquad=e^{\pm i\theta}
\bigl(k^\mu e^\nu(\mathbf k,\pm1)-k^\nu e^\mu(\mathbf k,\pm1)\bigr).
\tag{5.9.32}\label{eq:5.9.32}
\end{align}
This shows that the coefficient function that satisfies Eq.~\eqref{eq:5.9.6} for the
antisymmetric tensor representation of the homogeneous Lorentz group is (with
an appropriate choice of normalization)
\begin{equation}
u_{\mu\nu}(\mathbf p,\pm1)=\frac{i}{\sqrt{(2\pi)^3 2p^0}}
\bigl[p^\mu e^\nu(\mathbf p,\pm1)-p^\nu e^\mu(\mathbf p,\pm1)\bigr],
\tag{5.9.33}\label{eq:5.9.33}
\end{equation}
where $e^\mu(\mathbf p,\pm1)$ is given by Eq.~\eqref{eq:5.9.25}. Using this together
with Eq.~\eqref{eq:5.9.23} gives the general antisymmetric tensor field for massless
particles of helicity $\pm1$ in the form
\begin{equation}
f_{\mu\nu}=\partial_\mu a_\nu-\partial_\nu a_\mu.
\tag{5.9.34}\label{eq:5.9.34}
\end{equation}
Note that this is a tensor even though $a^\mu$ is not a four-vector, because
the extra term in Eq.~\eqref{eq:5.9.31} drops out in Eq.~\eqref{eq:5.9.34}. Note also that Eqs.
\eqref{eq:5.9.34}, \eqref{eq:5.9.24}, \eqref{eq:5.9.28}, and \eqref{eq:5.9.29} show that $f^{\mu\nu}$ satisfies the
vacuum Maxwell equations:
\begin{align}
\partial_\mu f^{\mu\nu}&=0,
\tag{5.9.35}\label{eq:5.9.35}\\
\epsilon^{\rho\sigma\mu\nu}\partial_\sigma f_{\mu\nu}&=0.
\tag{5.9.36}\label{eq:5.9.36}
\end{align}
To calculate the commutation relations for the tensor fields we need sums over
helicities of the bilinears $e^\mu e^{\nu*}$. The explicit formula \eqref{eq:5.9.21}
gives
\[
\sum_{\sigma=\pm1}e^i(\mathbf k,\sigma)e^j(\mathbf k,\sigma)^*
=\delta_{ij}-\frac{k^ik^j}{|\mathbf k|^2}
\]
and so, using Eq.~\eqref{eq:5.9.25},
\begin{equation}
\sum_{\sigma=\pm1}e^i(\mathbf p,\sigma)e^j(\mathbf p,\sigma)^*
=\delta_{ij}-\frac{p^ip^j}{|\mathbf p|^2}.
\tag{5.9.37}\label{eq:5.9.37}
\end{equation}
A straightforward calculation gives then
\begin{align}
[f_{\mu\nu}(x),f_{\rho\sigma}(y)^\dagger]
={}&(2\pi)^{-3}
[-\eta_{\mu\rho}\partial_\nu\partial_\sigma
+\eta_{\nu\rho}\partial_\mu\partial_\sigma
+\eta_{\mu\sigma}\partial_\nu\partial_\rho
-\eta_{\nu\sigma}\partial_\mu\partial_\rho]
\notag\\
&\cdot\int d^3p\,(2p^0)^{-1}
\bigl[|\kappa|^2e^{ip\cdot(x-y)}-|\lambda|^2e^{-ip\cdot(x-y)}\bigr].
\tag{5.9.38}\label{eq:5.9.38}
\end{align}
This clearly vanishes for $x^0=y^0$ if and only if
\begin{equation}
|\kappa|^2=|\lambda|^2
\tag{5.9.39}\label{eq:5.9.39}
\end{equation}
in which case since $f_{\mu\nu}$ is a tensor the commutator also vanishes for
all space-like separations. Eq.~\eqref{eq:5.9.39} also implies that the commutator of
the $a^\mu$ vanishes at equal times, and as we shall see in Chapter 8 this is
enough to yield a Lorentz-invariant $S$-matrix. The relative phase of the
creation and annihilation operators can be adjusted so that $\kappa=\lambda$;
the fields are then Hermitian if the particles are their own
charge-conjugates, as is the case for the photon.

Why should we want to use fields like $a^\mu(x)$ in constructing theories of
massless particles of spin one, rather than being content with fields like
$f^{\mu\nu}(x)$ with simple Lorentz transformation properties? The presence of
the derivatives in Eq.~\eqref{eq:5.9.34} means that an interaction density constructed
solely from $f_{\mu\nu}$ and its derivatives will have matrix elements that
vanish more rapidly for small massless particle energy and momentum than one
that uses the vector field $a_\mu$. Interactions in such a theory will have a
correspondingly rapid fall-off at large distances, faster than the usual
inverse-square law. This is perfectly possible, but gauge-invariant theories
that use vector fields for massless spin one particles represent a more
general class of theories, including those that are actually realized in
nature.

Parallel remarks apply to gravitons, massless particles of helicity $\pm2$.
From the annihilation and creation operators for such particles we can
construct a tensor $R_{\mu\nu\rho\sigma}$ with the algebraic properties of the
Riemann--Christoffel curvature tensor: antisymmetric within the pairs $\mu,\nu$
and $\rho,\sigma$, and symmetric between the pairs. However, in order to
incorporate the usual inverse-square gravitational interactions we need to
introduce a field $h_{\mu\nu}$ that transforms as a symmetric tensor, up to
gauge transformations of the sort associated in general relativity with
general coordinate transformations. Thus in order to construct a theory of
massless particles of helicity $\pm2$ that incorporates long-range
interactions, it is necessary for it to have a symmetry something like general
covariance. As in the case of electromagnetic gauge invariance, this is
achieved by coupling the field to a conserved `current' $\theta^{\mu\nu}$, now
with two spacetime indices, satisfying $\partial_\mu\theta^{\mu\nu}=0$. The
only such conserved tensor is the energy-momentum tensor, aside from possible
total derivative terms that do not affect the long-range behavior of the force
produced.\footnote{If $\theta^{\mu_1\cdots\mu_N}$ is a tensor current
satisfying $\partial_{\mu_1}\theta^{\mu_1\cdots\mu_N}=0$, then
$\int d^3x\,\theta^{0\mu_2\cdots\mu_N}$ is a conserved quantity that
transforms like a tensor of rank $N-1$. The only such conserved tensors are the
scalar `charges' associated with various continuous symmetries, and the
energy-momentum four-vector. The conservation of any other four-vector, or any
tensor of higher rank, would forbid all but forward collisions.} The fields of
massless particles of spin $j\geq3$ would have to couple to conserved tensors
with three or more spacetime indices, but aside from total derivatives there
are none, so \emph{high-spin massless particles cannot produce long-range
forces}.

\begin{center}
$*\ *\ *$
\end{center}

The problems we have encountered in constructing four-vector fields for
helicities $\pm1$ or symmetric tensor fields for helicity $\pm2$ are just
special cases of a more general limitation. To see this, let's consider how to
construct fields for massless particles belonging to arbitrary representations
of the homogeneous Lorentz group. As we saw in Section 5.6, any representation
$D(\Lambda)$ of the homogeneous Lorentz group can be decomposed into
$(2A+1)(2B+1)$-dimensional representations $(A,B)$, for which the generators
of the homogeneous Lorentz group are represented by
\[
(\mathcal J_{ij})_{a'b',ab}=\epsilon_{ijk}
[(J_k^{(A)})_{a'a}\delta_{b'b}+(J_k^{(B)})_{b'b}\delta_{a'a}],
\]
\[
(\mathcal J_{i0})_{a'b',ab}=-i
[(J_i^{(A)})_{a'a}\delta_{b'b}-(J_i^{(B)})_{b'b}\delta_{a'a}],
\]
where $\mathbf J^{(j)}$ are the angular-momentum matrices for spin $j$. For
$\theta$ infinitesimal, $D(R(\theta))=1+i\mathcal J_{23}\theta$, so Eqs.
\eqref{eq:5.9.12} and \eqref{eq:5.9.13} give
\[
\sigma u_{ab}(\mathbf k,\sigma)=(a+b)u_{ab}(\mathbf k,\sigma),
\]
\[
-\sigma v_{ab}(\mathbf k,\sigma)=(a+b)v_{ab}(\mathbf k,\sigma),
\]
and so $u_{ab}(\mathbf k,\sigma)$ and $v_{ab}(\mathbf k,\sigma)$ must vanish
unless $\sigma=a+b$ and $\sigma=-a-b$, respectively. Also, letting $\alpha$
and $\beta$ become infinitesimal in Eq.~\eqref{eq:5.9.14} gives
\begin{align*}
0&=(\mathcal J_{31}+\mathcal J_{01})_{ab,a'b'}
u_{a'b'}(\mathbf k,\sigma)\\
&=(J_2^{(A)}+iJ_1^{(A)})_{aa'}u_{a'b}(\mathbf k,\sigma)
+(J_2^{(B)}-iJ_1^{(B)})_{bb'}u_{ab'}(\mathbf k,\sigma),
\end{align*}
\begin{align*}
0&=(\mathcal J_{32}+\mathcal J_{02})_{ab,a'b'}
u_{a'b'}(\mathbf k,\sigma)\\
&=(-J_1^{(A)}+iJ_2^{(A)})_{aa'}u_{a'b}(\mathbf k,\sigma)
+(-J_1^{(B)}-iJ_2^{(B)})_{bb'}u_{ab'}(\mathbf k,\sigma),
\end{align*}
or more simply
\[
(J_1^{(A)}-iJ_2^{(A)})_{aa'}u_{a'b}(\mathbf k,\sigma)=0,
\]
\[
(J_1^{(B)}+iJ_2^{(B)})_{bb'}u_{ab'}(\mathbf k,\sigma)=0.
\]
These require that $u_{ab}(\mathbf k,\sigma)$ vanishes unless
\begin{equation}
a=-A,\qquad b=+B
\tag{5.9.40}\label{eq:5.9.40}
\end{equation}
and the same is obviously also true of $v_{ab}(\mathbf k,\sigma)$. Putting this
together, we see that a field of type $(A,B)$ can be formed only from the
annihilation operators for a massless particle of helicity $\sigma$ and the
creation operators for the antiparticle of helicity $-\sigma$, where
\begin{equation}
\sigma=B-A.
\tag{5.9.41}\label{eq:5.9.41}
\end{equation}

For instance, the $(\tfrac12,0)$ and $(0,\tfrac12)$ parts of the Dirac field
for a massless particle can only destroy particles of helicity $-\tfrac12$ and
$+\tfrac12$ respectively, and create antiparticles of helicity $+\tfrac12$ and
$-\tfrac12$, respectively. In the `two-component' theory of the neutrino,
there is only a $(\tfrac12,0)$ field and its adjoint, so neutrinos have helicity
$-\tfrac12$ and antineutrinos helicity $+\tfrac12$ in this theory.

By the same methods as in Section 5.7, it can be shown that the $(j,0)$ and
$(0,j)$ fields for massless particles of spin $j$ (i.e., helicity $\mp j$)
commute with each other and their adjoints at space-like separations if the
coefficients of the annihilation and creation terms in Eq.~\eqref{eq:5.9.1} satisfy Eq.
\eqref{eq:5.9.39}. The relative phase of the annihilation and creation operators may
then be adjusted so that these coefficients are equal. It is easy to see that
the fields for a massless particle of spin $j$ of type $(A,A+j)$ or $(B+j,B)$
are just the $2A$th or $2B$th derivatives of fields of type $(0,j)$ or $(j,0)$,
respectively, so these more general fields do not need to be considered
separately here.

We can now see why it was impossible to construct a vector field for massless
particles of helicity $\pm1$. A vector field transforms according to the
$(\tfrac12,\tfrac12)$ representation, and hence according to Eq.~\eqref{eq:5.9.19} can
only describe helicity zero. (It is, of course, possible to construct a vector
field for helicity zero---just take the derivative $\partial_\mu\phi$ of a
massless scalar field $\phi$.) The simplest covariant massless field for
helicity $\pm1$ has the Lorentz transformation type $(1,0)\oplus(0,1)$; that
is, it is an antisymmetric tensor $f_{\mu\nu}$. Similarly, the simplest
covariant massless field for helicity $\pm2$ has the Lorentz transformation
type $(2,0)\oplus(0,2)$: a fourth rank tensor which like the
Riemann--Christoffel curvature tensor is antisymmetric within each pair of
indices and symmetric between the two pairs.

The discussion of the inversions $P,C,T$ given in the previous section can be
carried over to the case of zero mass with only obvious modifications.
